% Original-source intake: Carnap 1950, first edition; distinct Zabell HHL10 chapter.
\chapter{Carnap: logical probability and inductive inference}
\label{sec:p15-carnap-logical-probability-and-induction}

\section{Two probabilities and a task of explication}
\label{sec:p15-carnap-explication}

Carnap's \emph{Logical Foundations of Probability} (1950, first edition)
does not identify induction with the assertion that a sufficiently long
list of instances entails a universal law. Its proposed inductive logic
determines a degree of confirmation \(c(h,e)\) of a hypothesis sentence
\(h\) relative to an evidence sentence \(e\). The hypothesis may be a
singular prediction, a statistical statement, or a universal law. A
calculation about their relation is distinct from the empirical work of
establishing the evidence and from deciding how to act
\cite[preface, pp.~v--vi; secs.~1--3, 45]{Carnap1950Probability}.

\begin{hbdefinition}{def:p15-carnap-two-probabilities}{Explicanda and explicata}
Carnap's \emph{probability\(_1\)} is the degree to which evidence supports
a hypothesis; \emph{probability\(_2\)} is a relative-frequency concept.
The first relates sentences, the second concerns a distribution of
properties in a population or sequence. Their mathematical analogies do
not make them the same concept. An \emph{explicandum} is a prior,
possibly vague concept to be clarified; an \emph{explicatum} is its exact
replacement in a constructed system. Similarity to the intended use,
exactness, fruitfulness, and simplicity are criteria for this replacement,
not a theorem that there is exactly one acceptable replacement.
\end{hbdefinition}

Once a language and a measure have been specified, a statement such as
\(c(h,e)=3/4\) can be determined without learning whether \(h\) or
\(e\) is actually true. This is the sense in which Carnap regards the
theorems of inductive logic as analytic. It does not turn the observation
report \(e\) into an analytic truth, prove a uniformity of nature, or
guarantee that the prescribed expectations succeed in the actual world.
His 1950 argument for the program and a mathematical proof inside the
program are different kinds of justification.

\section{A language, its states, and its structures}
\label{sec:p15-carnap-language}

\begin{hbdefinition}{def:p15-carnap-state-range}{Finite language and range}
Fix a finite domain of \(N\geq1\) distinct individuals named
\(a_1,\ldots,a_N\). The names designate different individuals and
exhaust this domain. Begin with \(q\geq1\) primitive unary predicates
\(P_1,\ldots,P_q\), interpreted as logically independent qualitative
properties. A \emph{state description} specifies exactly one of
\(P_j(a_i)\) and \(\neg P_j(a_i)\) for each pair \((i,j)\).
It may be written as a conjunction in a fixed order. Let \(\Omega_N\)
be the set of these complete descriptions. Boolean connectives have their
usual truth conditions; quantifiers range over the named domain; identity
is identity of its distinct members. For any sentence \(h\), its
\emph{range} is
\[
 R_N(h)\defeq\{z\in\Omega_N:z\models h\}.
\]
Logical truth, falsity, implication, and equivalence in this constructed
finite language mean, respectively, universal range, empty range, range
inclusion, and equality of ranges.
\end{hbdefinition}

Carnap's general finite languages also allow primitive relations: a
state then decides every ground instance of every primitive relation,
including instances with repeated arguments. Nothing in the definition
of an individual-renaming orbit below requires monadic predicates.
However, the closed quantitative formulas developed here concern the
monadic case. This is also the restriction of the greater part of the
1950 appendix's inferential calculations, not a claim that the general
theory of relations has been supplied
\cite[secs.~18, 26--27, 31; sec.~110A, p.~566]{Carnap1950Probability}.

\begin{hbdefinition}{def:p15-carnap-structure-width}{Structure descriptions and logical width}
Permuting the individual names acts on state descriptions. Two states
are \emph{isomorphic} if one is obtained from the other by such a
permutation. A \emph{structure description} is the disjunction of all
states in one orbit, in a fixed order. For languages with relations this
describes a complete relational structure up to individual renaming;
it is not in general just a list of predicate frequencies.

In the monadic language there are \(k\defeq2^q\) mutually exclusive
and exhaustive \(Q\)-properties: all conjunctions choosing \(P_j\)
or \(\neg P_j\) for each \(j\). A state is equivalently a word of
length \(N\) over these \(k\) properties. Its structure is specified
by the counts \((n_1,\ldots,n_k)\), with nonnegative integer entries
and sum \(N\). A molecular property \(M\) is a union of
\(Q\)-properties; its \emph{logical width} \(w(M)\) is their number.
Its relative width is \(w(M)/k\), not its actual population frequency.
An exhaustive disjoint division \(M_1,\ldots,M_p\) into nonempty
molecular properties has widths \(w_1,\ldots,w_p\) summing to \(k\).
\end{hbdefinition}

\begin{hbtheorem}{thm:p15-carnap-language-counts}{proved=annex:p15.carnap-language-counts}{States, structures, and widths}
In the monadic language, the \(Q\)-properties partition the possible
qualitative profiles. Every molecular property is uniquely a union of
them, and its width is therefore well defined. There are \(k^N\) states
and
\[
 r_N=\binom{N+k-1}{k-1}
\]
structure descriptions. The orbit with counts \((n_1,\ldots,n_k)\)
contains \(N!/\prod_i n_i!\) states. A conjunction specifying
\(d\) distinct primitive predicates or their negations has width
\(2^{q-d}\), for \(0\leq d\leq q\).
\end{hbtheorem}

Logical independence is substantive. If a primitive predicate means
\emph{raven} and another means \emph{black}, a contingent law relating
them is permissible, but a meaning stipulation making raven logically
imply black invalidates the unrestricted state space. Carnap also asks
whether the primitive predicates are qualitatively complete for the
intended universe. In 1950 he explicitly leaves open whether full
completeness is necessary or a weaker requirement suffices
\cite[sec.~18B, pp.~72--76]{Carnap1950Probability}. This is a
condition on application, not something established by counting states.

\section{Regular functions do not yet choose an inductive method}
\label{sec:p15-carnap-regular}

\begin{hbdefinition}{def:p15-carnap-regular-confirmation}{A regular measure and its confirmation function}
A \emph{regular} finite measure assigns each state a strictly positive
mass \(m_N(z)\), with \(\sum_{z\in\Omega_N}m_N(z)=1\).
Extend it to sentences by
\[
 m_N(h)\defeq\sum_{z\in R_N(h)}m_N(z),\qquad
 c_N(h,e)\defeq\frac{m_N(h\land e)}{m_N(e)}.
\]
The second definition is used only if \(R_N(e)\ne\varnothing\).
Regularity makes precisely these denominators positive. No value is
assigned here to confirmation on logically contradictory evidence.
A regular measure is \emph{symmetric in individuals} if isomorphic
states receive equal mass.
\end{hbdefinition}

These are Carnap's D55-1--4 and D90-1--2, expressed in the notation of
finite probability already developed in
Chapter~\ref{sec:p15-finite-probability}. The conditioning arithmetic is
not a new probability calculus. The distinctive problem is the
language-relative choice of \(m\), and thus of the confirmation function
\cite[secs.~55, 59, 90--91]{Carnap1950Probability}.

\begin{hbtheorem}{thm:p15-carnap-regular-laws}{proved=annex:p15.carnap-regular-laws}{Confirmation identities and renaming invariance}
For regular finite \(m_N\), confirmation lies in \([0,1]\), respects
logical equivalence of either argument, equals one when \(e\) implies
\(h\), and equals zero when \(e\) implies \(\neg h\).
It obeys complement and union identities and, when the displayed
conditioning events are nonempty,
\[
 c_N(h\land j,e)=c_N(h,e)c_N(j,e\land h).
\]
If \(m_N\) is symmetric and \(\pi\) is one permutation of the
individual names, then
\[
 m_N(\pi h)=m_N(h),\qquad
 c_N(\pi h,\pi e)=c_N(h,e).
\]
The same permutation must be applied to both arguments. Symmetry alone
does not determine the numerical value for a new individual's property.
\end{hbtheorem}

The last qualification matters: \(c(Pc,Pa\land Pb)=c(Pd,Pa\land Pb)\)
under renaming does not specify either value. It says neither that the
observations are independent nor that all structures have equal weight.
Also, a high \(c(h,e)\) and an \emph{increase} in confirmation are
different notions: additional evidence \(j\) is positively relevant to
\(h\) on \(e\) precisely when \(c(h,e\land j)>c(h,e)\), provided
both values are defined.

\section{The actual 1950 choice: equal structure weights}
\label{sec:p15-carnap-star}

\begin{hbdefinition}{def:p15-carnap-star}{The functions \(m^*\) and \(c^*\)}
Let the finite state space have \(r_N\) individual-renaming orbits.
Give each orbit total mass \(1/r_N\), and distribute that mass equally
among its states. Thus, if \(O_z\) is the orbit containing \(z\),
\[
 m_N^*(z)\defeq\frac1{r_N|O_z|},\qquad
 c_N^*(h,e)\defeq\frac{m_N^*(h\land e)}{m_N^*(e)}.
\]
For comparison let \(m_N^\dagger(z)\defeq1/|\Omega_N|\), with
confirmation \(c_N^\dagger\). The dagger labels the equal-state
alternative, not another name for \(c^*\).
\end{hbdefinition}

\begin{hbtheorem}{thm:p15-carnap-star-construction}{proved=annex:p15.carnap-star-construction}{Uniqueness, fitting finite systems, and the equal-state contrast}
The orbit construction is the unique normalized symmetric measure giving
all structure descriptions the same mass. It is regular. In the monadic
language its mass at a state with counts \((n_1,\ldots,n_k)\) is
\[
 m_N^*(z)=\frac{(k-1)!\prod_i n_i!}{(N+k-1)!}.
\]
Keeping the primitive predicates fixed, these measures fit under deleting
individuals: the marginal of \(m_{N+1}^*\) on the first \(N\)
individuals is \(m_N^*\). Consequently a sentence using only a fixed
finite list of individuals has the same mass in every larger finite
system. Under \(m_N^\dagger\), by contrast, an unmentioned
individual's profile is uniform and independent of the observed profiles.
For a molecular property of width \(w\), its confirmation on any
consistent evidence concerning only other individuals is \(w/k\).
\end{hbtheorem}

\begin{hbexample}{def:p15-carnap-three-states}{Three individuals and one property}
Take \(q=1\), \(k=2\), \(N=3\). The four structures have
\(0,1,2,3\) positive instances and each has mass \(1/4\) under
\(m^*\). The two constant states have mass \(1/4\) each; each of
the other six states has mass \(1/12\). With
\(e=P(a_1)\land P(a_2)\),
\[
 m^*(e)=\frac14+\frac1{12}=\frac13,\qquad
 c^*(P(a_3),e)=\frac{1/4}{1/3}=\frac34.
\]
Equal state weights instead give \(c^\dagger(P(a_3),e)=1/2\).
Both measures are symmetric in individuals. Only the first calculation
learns from the two observations in this example.
\end{hbexample}

This is the choice in the \emph{1950} appendix, sec.~110A,
pp.~562--566. Carnap explicitly does \emph{not} regard indifference
among structures as a compelling justification: structures have different
logical features. Nor does he regard simplicity as sufficient. His
proposal is to compare consequences with the intended inductive practice,
and he permits improvement or replacement of \(c^*\). The promised
systematic comparison with other methods is deferred there to a second
volume; it must not be reported as a proof already contained in 1950.

\section{Five forms of inference, with calculations}
\label{sec:p15-carnap-methods}

\begin{hbtheorem}{thm:p15-carnap-predictive}{proved=annex:p15.carnap-predictive}{Prediction, grouped properties, and estimation}
For a division with widths \(w_1,\ldots,w_p>0\), let \(e\) report
the categories of \(s\) specified individuals, with counts
\(s_1,\ldots,s_p\). For a disjoint further list of \(t\)
individuals, let \(h_{\boldsymbol t}\) report only its category counts
\(t_1+\cdots+t_p=t\). Assume \(N\geq s+t\). Then
\[
 c_N^*(h_{\boldsymbol t},e)
 =\frac{\prod_{j=1}^p\binom{s_j+t_j+w_j-1}{t_j}}
        {\binom{s+t+k-1}{t}}.
\]
The value is unchanged if \(e\) reports only its counts. In particular,
for a new individual and a molecular property \(M\) of width \(w\),
\[
 c_N^*(M(a),e)=\frac{s_M+w}{s+k}.
\]
Here \(e\) may report the full \(Q\)-profiles or just membership in
the chosen division. For \(s>0\), this last value is
\[
 \frac{s}{s+k}\frac{s_M}{s}+\frac{k}{s+k}\frac wk.
\]
It lies between observed frequency and relative width and approaches the
observed frequency along samples on which that frequency has a limit.
Its value is also the \(c^*\)-mean estimate of the relative frequency
of \(M\) in a nonempty unobserved list. If \(S_M\) is the number
of \(M\)-individuals in the whole domain, then
\[
 \mathbb E^*[S_M/N\mid e]
 =\frac{s_M+(N-s)(s_M+w)/(s+k)}{N}.
\]
\end{hbtheorem}

The block formula and its singular case are 1950 sec.~110C,
pp.~567--569, equations (6)--(7). A subsequent sample covering all
\(N-s\) remaining individuals gives the \emph{inverse} distribution
for the population counts: replace \(t_j\) by \(n_j-s_j\) and
\(t\) by \(N-s\), assigning zero to impossible counts. Thus inverse
inference is not a separate instruction to equate an unknown population
frequency to the sample frequency.

\begin{hbexample}{def:p15-carnap-width-calculation}{Widths change the rule of succession}
With two independent primitive properties \(P,R\), there are four
\(Q\)-properties. The widths of \(P\), \(P\land R\), and
\(P\lor R\) are respectively \(2,1,3\). Suppose ten complete
profiles yield counts \((6,2,1,1)\) in the order
\((PR,P\neg R,\neg PR,\neg P\neg R)\). Then
\[
 c^*(P(a),e)=\frac{8+2}{10+4}=\frac57,\quad
 c^*((P\land R)(a),e)=\frac{6+1}{14}=\frac12,\quad
 c^*((P\lor R)(a),e)=\frac{9+3}{14}=\frac67.
\]
If a new two-person sample is described only as one \(P\)-individual
and one non-\(P\)-individual, its confirmation is
\(\binom{10}{1}\binom4{1}/\binom{15}{2}=8/21\).
Applying \((s_M+1)/(s+2)\) indiscriminately to every molecular property
would not respect these widths or the relations between properties.
\end{hbexample}

\begin{hbtheorem}{thm:p15-carnap-direct}{proved=annex:p15.carnap-direct}{Direct inference from known population counts}
Let evidence specify only that the domain has category counts
\(n_1,\ldots,n_p\), summing to \(N\). For a fixed sample of
\(s\leq N\) individuals, the confirmation of its category-count
report \((u_1,\ldots,u_p)\) is, under \emph{any} regular symmetric
measure,
\[
 \frac{\prod_j\binom{n_j}{u_j}}{\binom Ns},\qquad
 \sum_j u_j=s,
\]
with value zero for impossible counts. The confirmation that one named
member belongs to category \(j\) is \(n_j/N\). The estimate of
its sample count is \(sn_j/N\). For fixed \(s\), as
\(N\to\infty\) and \(n_j/N\to r_j\), the count probability
converges to \(s!\prod_j r_j^{u_j}/\prod_j u_j!\).
\end{hbtheorem}

Carnap calls this \emph{direct} inference, from a population to a sample,
and distinguishes it from prediction between disjoint samples. Its
finite law is hypergeometric; the binomial or multinomial law is a
limit or approximation in this setting, not an exact law for sampling
without replacement from an arbitrary finite population
\cite[secs.~94--95, pp.~492--498]{Carnap1950Probability}.
For example, among five individuals exactly three are \(P\); the
confirmation that two specified individuals are both \(P\) is
\(\binom32/\binom52=3/10\), not \((3/5)^2\).

\begin{hbtheorem}{thm:p15-carnap-analogy}{proved=annex:p15.carnap-analogy}{The 1950 width-based inference by analogy}
Let \(B\subsetneq A\) be nonempty molecular properties with widths
\(0<w_B<w_A\), and let \(b,c\) be distinct individuals. With
\(e=B(b)\), \(j=A(c)\), and \(h=B(c)\),
\[
 c^*(h,j)=\frac{w_B}{w_A},\qquad
 c^*(h,e\land j)=\frac{w_B+1}{w_A+1}.
\]
The increase is \((w_A-w_B)/(w_A(w_A+1))>0\). In particular,
\(A=P\), \(B=P\land R\) in the two-predicate language gives
an increase from \(1/2\) to \(2/3\).
\end{hbtheorem}

This is sec.~110D, p.~569, equation (8), not a general metric theory of
similarity. It weighs the conjunction of shared properties through its
width. Counting verbal rephrasings of one shared property as separate
similarities cannot change that width.

\begin{hbtheorem}{thm:p15-carnap-universal}{proved=annex:p15.carnap-universal}{A law and its next instance have different confirmations}
Let \(M\) have width \(0<w<k\), and let evidence say that \(s\)
specified individuals all satisfy \(M\). Put
\(x^{\overline t}\defeq\prod_{i=0}^{t-1}(x+i)\), with empty
product one. In a domain of size \(N\geq s\),
\[
 c_N^*(\forall x\,M(x),e)
 =\frac{(s+w)^{\overline{N-s}}}{(s+k)^{\overline{N-s}}},\qquad
 c_N^*(M(a),e)=\frac{s+w}{s+k}\quad(N>s).
\]
For fixed \(s\), the first value tends to zero as \(N\to\infty\).
In the one-predicate case it is \((s+1)/(N+1)\), while the next
instance has value \((s+1)/(s+2)\).

More generally let \(U=\forall x(A(x)\to B(x))\). If the sample
contains \(s_-\) instances of \(A\land\neg B\) and \(s_+\)
of \(A\land B\), whose widths are \(w_-,w_+>0\), then its
unqualified next-instance confirmation and its qualified-instance
confirmation are respectively
\[
 1-\frac{s_-+w_-}{s+k},\qquad
 \frac{s_++w_+}{s_-+s_++w_-+w_+}.
\]
The second conditions additionally on the new individual's being \(A\).
\end{hbtheorem}

These are the distinction and calculations of sec.~110F--H,
pp.~570--575. After one hundred \(P\)-instances and no failures,
the next-instance confirmation in the one-predicate language is
\(101/102\); confirmation of \(\forall xP(x)\) in a population
of one million is \(101/1{,}000{,}001\). With two primitive
predicates interpreted as Swan and White, one hundred white swans and
no nonwhite swans give qualified-instance confirmation \(101/102\).
Non-swans contribute to the unqualified conditional instance but not
to that qualified ratio. A negative instance makes the original
universal law incompatible with evidence; it must not be quietly removed
by changing the hypothesis to a law about only unobserved individuals.

Carnap's interpretation of scientific reliance on a law is an argument:
an engineer needs sufficiently confirmed predictions about the bridges
being built, not a bet that no counterinstance occurs anywhere in an
unbounded universe. The predictions can be calculated directly from
observations. A high confirmation of the universal law is not a necessary
intermediate premise. This does not prove that laws have no scientific
role; Carnap explicitly retains their organizing and predictive usefulness.

\section{Application, total evidence, and objections}
\label{sec:p15-carnap-application}

\begin{hbmethod}{def:p15-carnap-application-method}{Applying an inductive calculation}
Specify the domain and primitive qualitative vocabulary; check the
independence assumptions behind the possible profiles. Translate the
actual hypothesis and evidence into this language, recording which
individuals are named and whether a report specifies identities or only
counts. Choose the measure and state its symmetry and additional
constraints. Check consistency of evidence before division. Calculate
the desired confirmation or estimate, identifying whether it is direct,
predictive, inverse, analogical, or universal inference. Finally, check
that omitted knowledge is irrelevant, or justify an approximation rather
than presenting a reduced evidence sentence as all available knowledge.
\end{hbmethod}

Carnap's \emph{total-evidence requirement} belongs to methodology: a
calculation on \(e\) licenses the corresponding belief only when \(e\)
expresses the available observational knowledge, or omitted knowledge
\(j\) is irrelevant in the required sense. Unlike deductive entailment,
the numerical value may rise \emph{or fall} when evidence is added.
In sec.~45B, pp.~211--213, he criticizes a sunrise calculation based
only on past sunrises for omitting knowledge relevant through mechanics;
he does not say probability is inapplicable to causally explained events.
Section~45C--D, pp.~213--218, argues for approximation by simplified
languages while warning that formal elegance may conceal omitted factors.

Three different objections must remain separate. First, symmetry of
individuals does not choose \(c^*\). Second, even after choosing a
measure, interpretation and total evidence can obstruct application.
Third, a successful exact calculation inside the system does not by itself
justify the system's rationality criteria. Section~110J, pp.~576--577,
distinguishes reconstructing an existing practice from vindicating its
truth or validity. The mathematical results above accomplish neither a
deductive answer to Hume nor a language-independent solution to Goodman's
projectibility problem.

\begin{hbexample}{def:p15-carnap-language-dependence}{A change of primitive vocabulary changes \(c^*\)}
In a language with only \(P\), three observed \(P\)-individuals give
\(4/5\) for a new \(P\)-individual. In a language with primitive
\(P,R\), the same report about \(P\) gives \((3+2)/(3+4)=5/7\).
The relative width of \(P\) is still \(1/2\), but the total prior
weight implicit in \(c^*\) has changed. This is not a failure of the
fitting theorem, which keeps the primitive vocabulary fixed and extends
only the individual domain. Nor may an actually known distinction be
discarded solely to obtain the more attractive number.
\end{hbexample}

\section{The historical chapter: development rather than replacement}
\label{sec:p15-carnap-history}

Zabell's distinct chapter, \emph{Carnap and the Logic of Inductive
Inference}, reads the history from a modern Bayesian perspective and
deliberately replaces much of Carnap's sentence notation by event and
probability notation. It contrasts an early search for one confirmation
function with later constrained families. That interpretive framework
must not be read back as the literal 1950 presentation
\cite[secs.~1--2, 13]{Zabell2011CarnapInductiveInference}.

Sections~3--5 trace the rule of succession through Bayes and Laplace to
Johnson's permutation and combination postulates. Individual symmetry
becomes exchangeability of the observation sequence; equal monadic
structure weights become equal probabilities for count vectors.
Johnson's later sufficientness postulate restricts a next-category
probability to the category's own count and the total sample size.
The historical characterization requires at least three categories;
with two, that sufficientness condition alone is vacuous. Zabell also
distinguishes Carnap's initial assumption of a predictive form from later
equivalence proofs by Kemeny and subsequent asymmetric versions.

\begin{hbdefinition}{def:p15-carnap-history-family}{A related later predictive family}
For \(k\) categories choose \(\alpha_i>0\) and put
\(\alpha=\sum_i\alpha_i\). Assign a specified finite word with
counts \(n_i\), length \(n\), the mass
\[
 m_{\boldsymbol\alpha}(z)\defeq
 \frac{\prod_i\alpha_i^{\overline{n_i}}}{\alpha^{\overline n}}.
\]
This is a related mathematical construction for the later discussion,
not a replacement for the 1950 choice. On complete \(Q\)-profiles,
\(\alpha_i=1\) gives the already defined \(m^*\).
\end{hbdefinition}

\begin{hbtheorem}{thm:p15-carnap-history-family}{proved=annex:p15.carnap-history-family}{A coherent family, not a characterization by symmetry alone}
The displayed masses are positive, normalized, exchangeable, and fit
under deleting the last observation. Their predictive probabilities are
\[
 \frac{n_i+\alpha_i}{n+\alpha}
 =\frac n{n+\alpha}\frac{n_i}{n}
  +\frac\alpha{n+\alpha}\frac{\alpha_i}{\alpha}\quad(n>0).
\]
Neither exchangeability nor exchangeability together with category
symmetry forces this family. For example, mix equally the six die
models in which one distinguished face has probability \(2/7\) and
the other five have probability \(1/7\). This mixture is exchangeable
and category-symmetric, but its predictive probability for face one is
\[
 \frac17\left(1+\frac{2^{n_1}}{\sum_{j=1}^6 2^{n_j}}\right),
\]
which depends on other category counts, not merely \((n_1,n)\).
\end{hbtheorem}

The construction and the obstruction clarify the claim of secs.~5--6:
the sufficientness postulate has content beyond symmetry. A convenient
Dirichlet representation does not make that postulate logically
compulsory. Zabell's example uses the two count vectors
\((20,10,10,10,10,10)\) and \((20,30,5,5,5,5)\); they have the
same first count and total but give different predictions for face one.

\begin{hbtheorem}{thm:p15-carnap-history-finite}{proved=annex:p15.carnap-history-finite}{Finite exchangeability, extendability, and instantial relevance}
Every exchangeable law on a finite category-word space is a mixture of
uniform laws on its count-vector orbits, with its count probabilities as
mixing weights. An exchangeable two-observation law assigning \(1/2\)
to each of \(RB,BR\) and zero to \(RR,BB\) has no exchangeable
extension to three observations. More generally, exchangeability of
\(m\) nondegenerate Bernoulli indicators gives common correlation
\(\rho\geq-1/(m-1)\). If such a law fits exchangeable extensions
of arbitrarily large length, \(\rho\geq0\), and
\(P(X_2=1\mid X_1=1)\geq P(X_2=1)\).
\end{hbtheorem}

These are the finite representation, extendability obstruction, and
variance argument of Zabell's secs.~7--8, pp.~281--288. Sampling
without replacement from a finite urn is a perfectly legitimate setting,
not pathological because it is not infinitely extendable. A known
three-red/two-black urn gives \(3/5\) initially but \(1/2\) for a
second red after a first red. Infinite extendability, not finite symmetry
alone, supports the nonnegative-relevance conclusion.

\begin{hbtheorem}{thm:p15-carnap-history-universal}{proved=annex:p15.carnap-history-universal}{A positive-law alternative and the first induction theorem}
In the binary case mix \(m^*\) with weight \(1-a\) and the
all-positive state law with weight \(a\), where \(0<a<1\).
After \(s\) positive observations, the finite-system confirmation of
the all-positive universal sentence has limit
\[
 \frac{a}{a+(1-a)/(s+1)}>0
\]
as the domain grows. This is a different measure, not a change of the
original \(c^*\) value.

Separately, let \(P(H)>0\) and let \(E_j\) satisfy
\(P(E_j\mid H)=1\) for every \(j\geq1\). Then
\[
 P(E_{n+1}\land\cdots\land E_{n+t}\mid E_1\land\cdots\land E_n)
 \longrightarrow1
\]
uniformly in finite \(t\geq1\). Exchangeability is not required;
the positive prior support for \(H\) is essential to this proof.
\end{hbtheorem}

The first calculation supplies a transparent example of the alternative
discussed in Zabell's sec.~6: zero confirmation of an infinite universal
law is not compulsory for all probability measures. It changes the
initial allocation. The second result is the first induction theorem
presented in his sec.~9, p.~288; its positive-hypothesis premise prevents
using it to overturn the original \(c^*\) universal-law limit.

The remainder of the historical account has a different scope from the
1950 appendix. Sections~10--11 follow analogy by proximity into partial
and Markov exchangeability, analogy by similarity into parametric
Bayesian models, and unknown categories into exchangeable partitions and
species-sampling laws. These are later developments, not the two-person
width calculation already proved here. Section~12 places conditional
probability puzzles, ravens, Goodman, total evidence, old evidence and
logical omniscience, and the Popper--Carnap controversy into this
historical setting; it does not reduce them all to a single mathematical
contradiction. We leave the independently allocated Hempel and Popper
presentations distinct.

Zabell's sec.~13 records substantive changes: movement from sentences
to events, acceptance of betting/coherence and subjective-probability
techniques, abandonment of the original predicate independence and
completeness requirements in later work, and a distinction between adult
credence and an idealized initial credibility. His objection is that
symmetry postulates are appropriate to particular epistemic situations,
not universally mandatory; an ideal initial language can relocate rather
than eliminate the choice of prior. His favorable historical judgment
about Carnap's influence on acceptance of epistemic probability is
compatible with that criticism of the strongest objectivist ambition.
These are attributed philosophical assessments, not theorems proving
the impossibility of every logical probability program.

\begin{proofboundary}{def:p15-carnap-source-boundary}{What this reconstruction establishes}
The annex supplies finite proofs of every numbered mathematical result
claimed above, including the grouped predictive rule, analogy,
law/instance distinction, and the bounded historical examples. They are
independent reconstructions, not a claim that the 1950 appendix printed
all those proofs. The original first-edition source, its later 1952
continuum and 1962 edition, and the separate historical chapter are not
interchangeable.

The comparative-confirmation, relevance and general estimation developments
continue in Chapter~\ref{sec:p15-carnap-confirmation-relevance-and-reliability}.
The infinite-system limits, the monadic existence and effective
reduction claims of sec.~110A with 110A (5), the sec.~41 interpretations,
the binomial law and Bernoulli's theorem, the sec.~110I requirement, and
a source-to-host bridge are given in Chapter~\ref{sec:p15-w02b-p15-029}.
The promised comparative vindication of \(c^*\) and a proof of the
variety-of-instances assertion, which the book refers to an earlier
paper, are not supplied. This account also does not
reconstruct all historical characterization and representation
proofs: Johnson--Carnap sufficientness, infinite de Finetti,
Hintikka--Niiniluoto extensions, the quantitative finite approximation,
Markov/normal mixture characterizations, and the general three-parameter
species-sampling theorem remain separate proof dependencies. The acquired
history copy has the allocated title and pp.~265--309, but a 2009
copyright footer; its identity with the publisher's 2011 chapter has not
been established by byte comparison. These limits prevent claiming
whole-source closure for either allocation.
\end{proofboundary}
